% Figure: multiplication by a unit complex number e^{i theta} as rotation.
\begin{figure}[htbp]
\centering
\begin{tikzpicture}[scale=1.6, line cap=round, >=latex]
  \draw[->, engblack, thin] (-1.6, 0) -- (2.0, 0) node[right, font=\small] {$\Re$};
  \draw[->, engblack, thin] (0, -1.6) -- (0, 2.0) node[above, font=\small] {$\Im$};
  % Reference unit circle
  \draw[engblack!40, thin, dashed] (0, 0) circle (1);
  \draw[engblack!40, thin, dashed] (0, 0) circle (1.414);
  % z = 1 + i (magnitude sqrt2, argument 45 deg)
  \draw[engblue, thick, ->] (0, 0) -- (1, 1);
  \filldraw[engblue] (1, 1) circle (0.04);
  \node[engblue, font=\footnotesize, anchor=south west] at (1, 1) {$z = 1 + i$};
  % Multiply by e^{i pi/2} = i: i * (1 + i) = -1 + i
  \draw[engred, thick, ->] (0, 0) -- (-1, 1);
  \filldraw[engred] (-1, 1) circle (0.04);
  \node[engred, font=\footnotesize, anchor=south east] at (-1, 1) {$i \cdot z = -1 + i$};
  % Arc of rotation
  \draw[engred, thick, ->] (1, 1) arc (45:135:1.414);
  \node[engred, font=\footnotesize] at (0, 1.65) {rotation by $\pi/2$};
  % Original argument arc
  \draw[engblue, thick] (0.35, 0) arc (0:45:0.35);
  \node[engblue, font=\footnotesize] at (0.55, 0.18) {$\pi/4$};
\end{tikzpicture}
\caption[Multiplication by a unit complex number is rotation]{The
unit complex number $e^{i\pi/2} = i$ multiplied into $z = 1 + i$
rotates $z$ counterclockwise through a quarter turn around the
origin, producing $iz = -1 + i$. The magnitude $\sqrt{2}$ is
preserved; only the argument changes. The same picture holds for
any unit complex number $e^{i\theta}$: multiplication by
$e^{i\theta}$ is rotation through angle $\theta$ around the
origin, with magnitude unchanged. The complex-multiplication rule
is the algebraic content of two-dimensional rotation.}
\label{fig:vol02:ch03:rotation}
\end{figure}
